\section{The Zero-Dimensional Limit: Recovery of Transformer Attention}\label{sec:transformers}

We now show that standard scaled dot-product attention is recovered as a gauge-fixed and isotropic-Gaussian limit of the construction; the full statement of what is and is not derived accompanies the boxed identity closing this section (fence F1; Section~\ref{sec:epistemic_status}). Standard attention is recovered along two complementary routes, whose reductions are collected as the hypotheses (R1)--(R3) and (N) of Theorem~\ref{thm:transformer_recovery} below. The first route (Section~\ref{sec:untied_qk}) retains $\Omega_{ij} = U_iU_j^{-1}$ non-trivially throughout and exposes untied query-key projections $Q_i = U_i^{-1}\mu_i$, $K_j = U_j^\top\Sigma_j^{-1}\mu_j$ directly from the transported KL, with the implicit pair-dependent bilinear ranging over $\mathrm{GL}(d_k)$ as the per-token frames and precisions vary; this route is structurally minimal in the sense that no object beyond the existing belief and frame fields is introduced, and it additionally captures the rotary positional structure of modern transformers. The second route (Section~\ref{sec:dot_product_derivation}) takes the reductions (R1)--(R3), yielding the familiar $QK^\top/\sqrt{d_k}$ form with the learned shared bilinear $M$ filling the cross-coupling slot vacated by the trivial transport; the precise overlap of the two routes is characterized after Eq.~\eqref{eq:gauge_bilinear}. Once a learned bilinear $M$ is introduced, factorization $M = W_QW_K^\top$ into query-key maps is purely algebraic, with existence and the exact non-uniqueness characterized in the discussion surrounding Eq.~\eqref{eq:head_space_kernel_participatory}. The reductions place standard attention inside the KL-attention family; whether trained transformers exploit the additional gauge structure is not tested here, the experiments of Section~\ref{sec:scaling_validation} train the general gauge architecture itself without these limits (so they constrain the gauge model directly rather than standard transformers), and behavioral claims about how humans coordinate meaning through linguistic tokens are outside scope (Section~\ref{sec:scope_limitations}).

\subsection{The Zero-Dimensional Limit and Gauge Fixing}

Consider the degenerate case where all agents exist at a single point $x_0 \in \mathcal{C}$ of the noumenal manifold, with $\mathrm{supp}(\sigma_i) = \{x_0\}$ for all agents $i$. The base manifold effectively becomes zero-dimensional, all spatial structure collapses, and agents interact purely through information-theoretic coupling at the single point. This is precisely the topology of standard transformers, where tokens (agents) are treated as an unordered set with no intrinsic spatial structure beyond sequence position.

\subsection{Setup: Agents as Gaussian Beliefs in Local Frames}

In the general formulation, each agent $i$ (token) maintains a local MVG belief
\begin{equation}
q_i = \mathcal{N}(\mu_i, \Sigma_i),
\end{equation}
where $\mu_i \in \mathbb{R}^{d_k}$ is the agent's mean representation in its local gauge frame, and $\Sigma_i \in \mathbb{R}^{d_k\times d_k}$ is its belief covariance (symmetric positive definite). We embed in $K_q = d_k$ dimensions to match standard notation, with gauge group the fundamental $\mathrm{GL}(d_k)$ until the multi-head restriction below. In modern multi-head transformer architectures the embedding dimension $d_{\text{model}}$ is partitioned into $H$ heads of dimension $d_{\text{head}} = d_{\text{model}}/H$ and each head carries its own gauge frame, so $K_q=d_k$ here corresponds to the per-head dimension $d_{\text{head}}$ rather than the full model width; the multi-head gauge group is the block-diagonal $\mathrm{GL}(d_{\text{head}})^H$ (Eq.~\eqref{eq:multihead_gauge_group}), identified in detail in the multi-head paragraph following the boxed attention formula.

Communication between agents $i$ and $j$ is mediated by a gauge frame transport operator
\begin{equation}
\Omega_{ij} \in \mathrm{GL}(d_k),
\end{equation}
which maps representations from agent $j$'s local frame into agent $i$'s frame. The value-aggregation output received by agent $i$ is the posterior mean of the mixture variational model under responsibilities $\beta_{ij}$ (derived formally in Section~\ref{sec:value_aggregation_itfb}):
\begin{equation}
\hat\mu_i = \sum_j \beta_{ij} \Omega_{ij}\mu_j,
\end{equation}
the gauge-theoretic analog of the attention aggregation $\sum_j \alpha_{ij} V_j$ in standard transformers.

\subsection{Connection to Standard Transformers}

The gauge-covariant attention $\beta_{ij} = \operatorname{softmax}(-D_{\mathrm{KL}}(q_i \| \Omega_{ij} q_j)/\tau)$ derived in Section~\ref{sec:mixture_derivation} is already a complete attention mechanism; standard transformer attention is recovered as a special case along the two routes stated in the section opening. The limits shared by both routes are developed first.

\paragraph{Deterministic beliefs via scaled limit.}

We consider a single-point base space $\mathcal{C}$ with a finite set of agents $\{1,\dots,N\}$ (tokens) with no particular ordering.

The naive deterministic limit fails: setting $\Sigma_i \to 0$ yields Dirac delta beliefs $q_i(k_i) \to \delta(k_i - \mu_i)$, but the KL divergence between distinct Diracs is infinite:
\begin{equation}
D_{\mathrm{KL}}(\delta(k - \mu_i) \| \delta(k - \mu_j)) = +\infty \quad \text{for } \mu_i \neq \mu_j.
\label{eq:dirac_kl}
\end{equation}
This divergence reflects the absolute continuity requirement in the definition of KL: $D_{\mathrm{KL}}(P \| Q)$ is finite only when $P$ is absolutely continuous with respect to $Q$ (i.e., $P(A) > 0$ implies $Q(A) > 0$ for all measurable $A$). Distinct Dirac measures have disjoint support and therefore violate this condition. We therefore take a joint scaling limit in which the belief variance $\sigma^2$ remains finite and is absorbed into learned parameters.

Next, we assume isotropic covariances $\Sigma_i = \Sigma_j = \sigma^2 I$ with $\sigma^2 > 0$ (hypothesis (R1) of Theorem~\ref{thm:transformer_recovery} below). This implies that the log-determinant terms cancel ($\log\frac{|\Sigma_j|}{|\Sigma_i|} = 0$) and the trace term simplifies to $\mathrm{Tr}(\Sigma_j^{-1}\Sigma_i) = \mathrm{Tr}(I) = d_k$.

With gauge transport $\Omega_{ij} \in \mathrm{GL}(d_k)$, the transported Gaussian $\Omega_{ij}q_j$ has mean $\Omega_{ij}\mu_j$ and covariance $\Omega_{ij}\Sigma_j\Omega_{ij}^\top = \sigma^2(\Omega_{ij}\Omega_{ij}^\top)$. Therefore the full KL divergence is
\begin{equation}
D_{\mathrm{KL}}(q_i \| \Omega_{ij}q_j) = \frac{1}{2}\Big[\log\det(\Omega_{ij}\Omega_{ij}^\top) + \mathrm{Tr}\big((\Omega_{ij}\Omega_{ij}^\top)^{-1}\big) - d_k\Big] + \frac{1}{2\sigma^2}(\mu_i - \Omega_{ij}\mu_j)^\top(\Omega_{ij}\Omega_{ij}^\top)^{-1}(\mu_i - \Omega_{ij}\mu_j).
\end{equation}

The Mahalanobis quadratic form simplifies via the identity $(\Omega_{ij}\Omega_{ij}^\top)^{-1} = \Omega_{ij}^{-\top}\Omega_{ij}^{-1}$:
\begin{equation}
(\mu_i - \Omega_{ij}\mu_j)^\top(\Omega_{ij}\Omega_{ij}^\top)^{-1}(\mu_i - \Omega_{ij}\mu_j) = \|\Omega_{ij}^{-1}(\mu_i - \Omega_{ij}\mu_j)\|^2 = \|\Omega_{ij}^{-1}\mu_i - \mu_j\|^2.
\label{eq:mahalanobis_identity}
\end{equation}

Rather than transporting the key belief into the query's frame via $\Omega_{ij}$ and measuring with the induced metric, the KL divergence equivalently transports the query into the key's frame via $\Omega_{ij}^{-1}$ and measures the Euclidean distance there. Either perspective is appropriate; the Mahalanobis metric automatically accounts for the distortion introduced by gauge transport. The KL alignment functional generalizes to the R\'enyi family $D_\alpha$ on its finiteness domain; the generalization is developed in the paragraph closing this section.

\paragraph{Isotropic attention with general gauge transport.}

The structural terms $\log\det(\Omega_{ij}\Omega_{ij}^\top)$, $\mathrm{Tr}((\Omega_{ij}\Omega_{ij}^\top)^{-1})$, and $d_k$ depend on the gauge transport but not on the belief means. Defining the geometric bias
\begin{equation}
S(\Omega_{ij}) := \frac{1}{2}\Big[\log\det(\Omega_{ij}\Omega_{ij}^\top) + \mathrm{Tr}\big((\Omega_{ij}\Omega_{ij}^\top)^{-1}\big) - d_k\Big],
\end{equation}
the KL divergence under isotropic beliefs with general transport takes the form
\begin{equation}
D_{\mathrm{KL}}(q_i \| \Omega_{ij} q_j) = S(\Omega_{ij}) + \frac{1}{2\sigma^2}\|\Omega_{ij}^{-1}\mu_i - \mu_j\|^2.
\end{equation}

The attention weight
\begin{equation}
\beta_{ij} = \operatorname{softmax}_j\left(-\frac{S(\Omega_{ij}) + \frac{1}{2\sigma^2}\|\Omega_{ij}^{-1}\mu_i - \mu_j\|^2}{\tau}\right),
\end{equation}
combines a content-dependent term (the Mahalanobis distance between belief means, measured in the key's frame) with a pair-dependent geometric bias $S(\Omega_{ij})$ that encodes the distortion introduced by gauge transport. When $\Omega_{ij}$ is orthogonal, $S(\Omega_{ij}) = 0$; for general $\mathrm{GL}(K_q)$ transport, $S(\Omega_{ij}) \geq 0$ with equality iff $\Omega_{ij} \in \mathrm{O}(K_q)$. The geometric bias therefore penalizes transport operators that distort volume or stretch directions, providing a geometric regularizer intrinsic to the KL divergence. This is the natural form of attention under a single simplifying limit (isotropic covariances); standard transformer attention requires a further specialization.

The rescaled KL divergence $\widetilde{D}_{\mathrm{KL}}(q_i \| \Omega_{ij}q_j) := \sigma^2 \cdot D_{\mathrm{KL}}(q_i \| \Omega_{ij}q_j) \to \frac{1}{2}\|\Omega_{ij}^{-1}\mu_i - \mu_j\|^2$ as $\sigma \to 0$ remains finite and equals half the squared distance between means as measured in the key's reference frame.

\paragraph{From transport to query-key compatibility.}

A constant-gauge presentation, in which a single nontrivial $\Omega$ is shared by all pairs $i,j$, is structurally inconsistent with the vertex-frame parameterization of this manuscript, as we now show. Under $\Omega_{ij} = U_iU_j^{-1}$ the self-pair gives $\Omega_{ii} = I$, so pair-independence $\Omega_{ij} = \Omega$ that includes the diagonal forces $\Omega = I$ directly; alternatively, restricting pair-independence to off-diagonal pairs, for $N \geq 3$ the cocycle $\Omega_{ij}\Omega_{jk} = \Omega_{ik}$ on three distinct indices yields $\Omega^2 = \Omega$, and invertibility forces $\Omega = I$. The only constant pair-independent transport consistent with the vertex-frame law is the trivial one. The reduction below therefore separates two distinct objects:
\begin{enumerate}
\item The \emph{relative gauge transport} $\Omega_{ij} = U_iU_j^{-1}$, which carries the gauge-covariant content of the framework. The transformer reduction sets $U_i = U$ for all tokens (a single shared frame), giving $\Omega_{ij} = I$ on the relevant base point.
\item A \emph{shared query-key bilinear compatibility} $M \in \mathrm{GL}(d_k)$ that encodes how query means are paired with key means at the level of the attention score. Standard transformer attention $\beta_{ij} \propto \exp\big(\mu_i^\top M \mu_j / \tau\big)$ corresponds to a learned bilinear form $M$ acting on the shared representation in the head's $d_k$-dimensional subspace, not to a non-trivial pair-independent gauge transport. In the rectangular multi-head setting with projections $W_Q, W_K \in \mathbb{R}^{d_{\text{model}} \times d_{\text{head}}}$, the bilinear $W_Q W_K^\top \in \mathbb{R}^{d_{\text{model}}\times d_{\text{model}}}$ is rank-deficient ($\mathrm{rank} \le d_{\text{head}}$) and is not an element of $\mathrm{GL}(d_{\text{model}})$; the invertible kernel $M_h$ lives in the per-head image subspace and is recovered via the thin-SVD decomposition $W_Q = U_Q A_Q$, $W_K = U_K A_K$ developed below in Eqs.~\eqref{eq:thin_svd_decomposition}--\eqref{eq:low_rank_lift}.
\end{enumerate}
With this separation, the geometric bias $S(\Omega_{ij}) = S(I) = 0$ is automatically constant and cancels under softmax; what remains is a Mahalanobis-style attention score whose bilinear part is the learned compatibility $M$. The cross term in the attention score takes the form $\mu_i^\top M \mu_j$, and the $\Omega \to I$ specialization removes any structural inconsistency between pairwise transport and pair-independent compatibility.

\subsection{Untied Query-Key Carving from Per-Token Frames}
\label{sec:untied_qk}

The trivial-frame reduction is one realization of standard attention. A second, complementary reduction retains $\Omega_{ij} = U_iU_j^{-1}$ throughout and exhibits standard attention with untied query-key projections directly from the gauge-covariant KL, without inserting a separate bilinear $M$ by hand.

Under arbitrary SPD covariances $\Sigma_i$ (independent of the frames $U_i$), the transported Gaussian KL admits the exact decomposition
\begin{multline}
D_{\mathrm{KL}}(q_i \| \Omega_{ij\#}q_j) = \tfrac{1}{2}\Big[\log\det\Sigma_j - \log\det\Sigma_i + 2\log|\det U_i| - 2\log|\det U_j| \\
 - d_k + \mathrm{tr}(H_j P_i) + x_i^\top H_j x_i - 2x_i^\top k_j + r_j\Big],
\label{eq:full_kl_general}
\end{multline}
with the common-frame quantities
\begin{equation}
x_i = U_i^{-1}\mu_i, \qquad P_i = U_i^{-1}\Sigma_i U_i^{-\top}, \qquad H_j = U_j^\top \Sigma_j^{-1} U_j, \qquad k_j = U_j^\top \Sigma_j^{-1}\mu_j, \qquad r_j = \mu_j^\top \Sigma_j^{-1}\mu_j.
\end{equation}
Equation~\eqref{eq:full_kl_general} is the inhomogeneous form taken by the gauge-covariant Gaussian KL when the covariance field is unconstrained and the frames are general $\mathrm{GL}(d_k)$ elements. It follows from the Gaussian KL with transported moments $(\Omega_{ij}\mu_j, \Omega_{ij}\Sigma_j\Omega_{ij}^\top)$ by direct substitution: the log-determinant block yields the four determinant terms, cyclicity of the trace yields $\mathrm{tr}(H_jP_i)$, and the quadratic block expands in the common-frame variables via $\Omega_{ij}^{-1}\mu_i = U_jx_i$; it is confirmed numerically to machine precision.

The cross term $-2x_i^\top k_j$ is the only coupling between query and key indices through the belief means. It carves cleanly as a dot product between per-token vectors,
\begin{equation}
x_i^\top k_j = (U_i^{-1}\mu_i)^\top (U_j^\top \Sigma_j^{-1}\mu_j) = Q_i^\top K_j,
\label{eq:untied_cross_term}
\end{equation}
where the gauge-derived query and key projections are
\begin{equation}
\boxed{Q_i = U_i^{-1}\mu_i, \qquad K_j = U_j^\top \Sigma_j^{-1}\mu_j.}
\label{eq:gauge_qk}
\end{equation}

For the same token, the two projections satisfy $W_Q^{(i)} = U_i^{-1}$ and $W_K^{(i)} = U_i^\top \Sigma_i^{-1}$, with the inverse-transpose tying broken by the precision factor $\Sigma_i^{-1}$. The construction is therefore genuinely untied: $W_Q$ and $W_K$ are different functions of the per-token belief, not the same projection up to symmetry.

The implicit pair-dependent bilinear is
\begin{equation}
M_{ij} := \Omega_{ij}^{-\top}\Sigma_j^{-1} = U_i^{-\top} U_j^\top \Sigma_j^{-1}.
\label{eq:gauge_bilinear}
\end{equation}

As $(U_i, U_j)$ ranges over $\mathrm{GL}(d_k)^2$ and $\Sigma_j$ over $\mathrm{SPD}(d_k)$, polar decomposition shows that the per-pair bilinear $M_{ij}$ ranges over all of $\mathrm{GL}(d_k)$, matching the pointwise range of an independently learned $W_Q W_K^\top \in \mathrm{GL}(d_k)$ in standard attention. The function classes nonetheless differ: standard attention learns one shared, token-independent bilinear, whereas $M_{ij}$ is pair-dependent and determined by the per-token frames and precisions. A shared token-independent $M$ is realizable from the gauge family only in the symmetric positive-definite case $M = \Sigma^{-1}$ (shared frames and a shared covariance), and the carving forces positive self-scores $Q_j^\top K_j = \mu_j^\top\Sigma_j^{-1}\mu_j > 0$, so the two function classes are incomparable: standard attention with $M = -I$ realizes negative self-scores, which the gauge family cannot, while the pair dependence of $M_{ij}$ is unavailable to any shared bilinear outside the symmetric positive-definite case $M = \Sigma^{-1}$; neither class contains the other. Generically $M_{ij}$ is non-symmetric, recovering the asymmetry of $W_QW_K^\top$ that this route's own symmetric $C^{-1}$ closure (next paragraph) gives up. Although $M_{ij}$ depends on the pair $(i,j)$, its evaluation $\mu_i^\top M_{ij}\mu_j$ factorizes through per-token vectors via Eqs.~\eqref{eq:untied_cross_term}--\eqref{eq:gauge_qk}, exactly as in standard attention.

\paragraph{Reduction to the standard form.}
The remaining terms in Eq.~\eqref{eq:full_kl_general} fall into two groups under the row-softmax over $j$. Terms depending on $j$ alone, namely $r_j$, $\log\det\Sigma_j$, and $\log|\det U_j|$, contribute a key-side prior bias that is absorbed into $\log\pi_{ij}$. Under the symmetric closure $\Sigma_j = U_j C U_j^\top$ the surviving $j$-dependent residual is exactly the key self-norm $r_j = \mu_j^\top\Sigma_j^{-1}\mu_j = \|K_j\|^2$ (with $K_j = C^{-1/2}U_j^{-1}\mu_j$), the same term the dot-product route of Section~\ref{sec:dot_product_derivation} carries; under the row softmax it cancels only when $\|K_j\|$ is constant in $j$ (key normalization), or when it is absorbed into a content-dependent per-key bias, in which case the recovered object is the key-bias-augmented transformer rather than the standard architecture. Treating $r_j$ as part of the static positional prior $\pi_{ij}$ is legitimate only to the extent it is read as such a bias, since $r_j$ is genuinely belief-dependent; so this route, like the dot-product route, recovers standard attention only up to the key-norm condition. Terms coupling $i$ and $j$ through the belief covariances rather than the means, namely $x_i^\top H_j x_i$ and $\mathrm{tr}(H_j P_i)$, are gauge-theoretic uncertainty corrections beyond vanilla attention; they vanish under the closure $\Sigma_j = U_j C U_j^\top$ for shared SPD $C$, which collapses $H_j$ to the constant $C^{-1}$. Under this closure the carving in Eq.~\eqref{eq:gauge_qk} becomes the symmetric form $Q_i = C^{-1/2}U_i^{-1}\mu_i$, $K_j = C^{-1/2}U_j^{-1}\mu_j$, and attention with a single shared symmetric bilinear $C^{-1}$ is recovered: the same closure that removes the covariance couplings also ties the projections, so the generic untied non-symmetric form of standard attention is not reached along this route.

\paragraph{Identification with rotary positional structure.}
The natural identification of the per-token frame $U_i$ with a real transformer architecture is the per-position rotational frame of rotary positional embeddings, in which $U_i \in \mathrm{O}(d_k)$ is a block-diagonal rotation depending on token position. In that case $U_i^{-1} = U_i^\top$, and the closure $\Sigma_j = U_j C U_j^\top$ reduces to the homogeneity condition $\Sigma_j = C$ for all $j$ (since $U_j C U_j^\top = C$ when $[C, U_j] = 0$): it is satisfied by any single shared covariance that is block-isotropic on the rotation planes, not by an arbitrary covariance field. Under that homogeneity assumption the carving in Eq.~\eqref{eq:gauge_qk} reduces, up to the shared factor $C^{-1/2}$, to $Q_i = U_i^\top\mu_i$ and $K_j = U_j^\top\mu_j$, which is the rotation-modulated query-key structure used in rotary attention, with $\mu_i$ playing the role of the position-independent content vector. This symmetric closure recovers only the rotational transport structure; actual rotary attention additionally applies distinct learned projections $W_Q \neq W_K$ before the rotations. The trivial-frame specialization $U_i = U$ for all $i$ removes the per-token positional variation; in that limit $\Omega_{ij} = I$ and the analysis of Section~\ref{sec:dot_product_derivation} below applies, with the learned bilinear $M$ replacing the structural role played by per-token frame variation.

\subsection{Recovery of Dot-Product Attention as a Gauge-Fixed Limit}
\label{sec:dot_product_derivation}

\begin{theorem}[Transformer recovery]
\label{thm:transformer_recovery}
Assume the following reductions of the gauge-covariant KL attention.
\begin{itemize}
\item[(R1)] Isotropic shared covariances: $\Sigma_i = \sigma^2 I$ for all $i$, with $\sigma^2 > 0$ held finite in the joint scaling limit above.
\item[(R2)] Shared frames: $U_i = U$ for all $i$, hence trivial pairwise transport $\Omega_{ij} = I$. The weaker hypothesis $\Omega_{ij} \in \mathrm{O}(d_k)$ removes only the geometric bias $S(\Omega_{ij})$ and does not suffice for the cross-term reduction below.
\item[(R3)] A separately learned shared bilinear $M = \sigma^2 W_QW_K^\top \in \mathrm{GL}(d_k)$ replacing the cross term, an object external to the gauge data.
\item[(N)] Key-norm constancy: $\|K_j\|$ is constant in $j$, or the residual key bias is absorbed into a learned per-key bias.
\end{itemize}
Then the attention weights reduce to $\beta_{ij} = \mathrm{softmax}_j\big(Q_iK_j^\top/(\kappa\sqrt{d_k})\big)$ with $Q_i = \mu_i^\top W_Q$ and $K_j = \mu_j^\top W_K$, and composing the value aggregation of Section~\ref{sec:value_aggregation_itfb} with a source-only value map $W_V$ recovers scaled dot-product attention at $\kappa = 1$. Hypothesis (R3) makes this a recovery, not a derivation (fence F1; Section~\ref{sec:epistemic_status}).
\end{theorem}

The paragraphs of this subsection constitute the proof: the expansion of the score under (R1)--(R3), the softmax cancellations, the key-bias analysis that (N) closes, the algebraic factorization that (R3) supplies, and the temperature scaling.

\paragraph{Expanding the KL Divergence.}
Under hypothesis (R2) the geometric bias collapses, $S(\Omega_{ij}) = 0$, and with (R1) the KL reduces to $\frac{1}{2\sigma^2}\|\mu_i - \mu_j\|^2$, whose cross term is $\mu_i^\top\mu_j$. Replacing that cross term by the learned bilinear $\mu_i^\top M \mu_j$ of hypothesis (R3) replaces the KL rather than reducing it: the resulting score is itself a divergence only for orthogonal $M$, although the row-constant discrepancy cancels under the softmax, leaving $\beta$ equal to KL-attention under a constant transport $\Omega = M^{-\top}$ of exactly the kind the vertex-frame law excludes. The substitution gives
\begin{equation}
s_{ij} = \frac{1}{2\sigma^2}\|\mu_i\|^2 + \frac{1}{2\sigma^2}\|\mu_j\|^2 - \frac{1}{\sigma^2}\mu_i^\top M \mu_j + C,
\end{equation}
where $M \in \mathrm{GL}(d_k)$ is the shared bilinear similarity in the head's $d_k$-dimensional subspace, with the identification $M = \sigma^2 W_Q W_K^\top$ holding for the square per-head projections $W_Q, W_K \in \mathbb{R}^{d_k \times d_k}$; the rectangular multi-head case $W_Q, W_K \in \mathbb{R}^{d_{\text{model}}\times d_{\text{head}}}$ is recovered via thin SVD on the head subspace, as discussed below. The constant $C$ is a $\mu$-independent normalization absorbed by the softmax.

\paragraph{Softmax cancellations.}
For a fixed query agent $i$, the softmax over keys $j$ is
\begin{equation}
\beta_{ij} = \frac{\exp(-s_{ij}/\tau)}{\sum_k \exp(-s_{ik}/\tau)}.
\end{equation}
The $i$-terms $\frac{1}{2\sigma^2}\|\mu_i\|^2 + C$ are independent of $j$ and cancel between numerator and denominator of the softmax, yielding
\begin{equation}
\beta_{ij} \propto \exp\left(\frac{1}{\tau}\left[\frac{1}{\sigma^2}\mu_i^\top M \mu_j - \frac{1}{2\sigma^2}\|\mu_j\|^2\right]\right).
\end{equation}
Thus, the effective logit is
\begin{equation}
\tilde{s}_{ij} = \frac{1}{\sigma^2}\mu_i^\top M \mu_j - \frac{1}{2\sigma^2}\|\mu_j\|^2.
\end{equation}

\paragraph{Key bias cancellation.}
The second term $-\frac{1}{2\sigma^2}\|\mu_j\|^2$ depends on key $j$ and is a key-dependent bias on the raw embedding norms. Two candidate mechanisms could cancel it; the first fails quantitatively, and only the second yields exact cancellation under stated conditions.

\textbf{(1) High-dimensional concentration does not suffice.}
For embeddings in $\mathbb{R}^{d_k}$ with approximately independent components of per-component variance $\sigma_0^2$ (e.g., $\mu_j^{(i)} \sim \mathcal{N}(0, \sigma_0^2)$), concentration of measure implies
\begin{equation}
\|\mu_j\|^2 = \sum_{i=1}^{d_k} (\mu_j^{(i)})^2 = d_k\sigma_0^2 \pm O(\sigma_0^2\sqrt{d_k}),
\end{equation}
with relative fluctuations $O(1/\sqrt{d_k})$. The softmax, however, already cancels the constant mean $d_k\sigma_0^2$ across keys; what matters is the across-key fluctuation of the bias, of absolute size $O(\sigma_0^2\sqrt{d_k})$, which is the same order as the spread of the content logits $\mu_i^\top M\mu_j$. After the temperature $\tau = \kappa\sqrt{d_k}$ both contribute $O(1)$ logit variation, so concentration alone never makes the bias negligible relative to the signal, at any $d_k$.

\textbf{(2) Norm normalization, with caveats.}
Modern transformer architectures apply layer normalization to the input vector $\mu_j$ before the learned key projection $K_j = W_K^\top \mu_j$. Layer normalization makes $\|\mu_j\|$ approximately constant across tokens but does \emph{not} by itself force $\|K_j\|$ to be constant: the projected norm depends on the spectrum of $W_K$ and is constant in $j$ only if $W_K$ is orthogonal/scaled or if a separate post-projection key normalization is applied. Hypothesis (N) is therefore realized in three exact regimes. First, $\|\mu_j\|$ exactly constant in $j$ (pre-projection layer normalization that projects embeddings to a common sphere) makes the bias row-constant, so it cancels under the softmax. Second, $\|K_j\|$ constant in $j$ (post-projection normalization, sometimes denoted query/key norm in modern variants): post-projection key normalization controls $\|K_j\|$, not $\|\mu_j\|$, so it cancels the bias exactly along the untied route of Section~\ref{sec:untied_qk}, where the residual key-side term is $r_j = \|K_j\|^2$, and along the present route only when the Gaussian beliefs are taken over the projected key coordinates, identifying $\mu_j$ with $K_j$ so that the bias becomes $\|K_j\|^2$. Third, the key bias $-\tfrac{1}{2\sigma^2}\|\mu_j\|^2$ may be absorbed into a learned per-key bias term, which describes a transformer augmented with such a bias rather than the standard architecture. Outside these regimes, the KL limit gives a distance/RBF-style attention with a residual key-dependent bias rather than pure dot-product attention.

\paragraph{Factorization into query-key dot product.}
The leading term in $\tilde{s}_{ij}$ is
\begin{equation}
\frac{1}{\sigma^2}\mu_i^\top M \mu_j.
\end{equation}
We factor the bilinear similarity $M$ in the head subspace as
\begin{equation}
W_Q W_K^\top = \frac{1}{\sigma^2} M \in \mathrm{GL}(d_k),
\label{eq:wqwk_square}
\end{equation}
which holds directly when $W_Q, W_K \in \mathbb{R}^{d_k \times d_k}$ are square per-head projections. For the standard rectangular multi-head case $W_Q, W_K \in \mathbb{R}^{d_{\text{model}}\times d_{\text{head}}}$ with $d_{\text{head}} = d_k < d_{\text{model}}$, each rectangular projection is decomposed via thin singular-value decomposition into an isometric subspace embedding and an invertible head-space map,
\begin{equation}
W_Q = U_Q A_Q, \qquad W_K = U_K A_K,
\label{eq:thin_svd_decomposition}
\end{equation}
where $U_Q, U_K \in \mathbb{R}^{d_{\text{model}}\times d_{\text{head}}}$ have orthonormal columns ($U_Q^\top U_Q = U_K^\top U_K = I_{d_{\text{head}}}$, isometric subspace embeddings) and $A_Q, A_K \in \mathrm{GL}(d_{\text{head}})$ are invertible head-space maps. The invertible head-space kernel is then
\begin{equation}
M_h := A_Q A_K^\top \in \mathrm{GL}(d_{\text{head}}),
\label{eq:head_space_kernel_participatory}
\end{equation}
and the ambient transformer logit kernel is the low-rank lift
\begin{equation}
W_Q W_K^\top = U_Q M_h U_K^\top \in \mathbb{R}^{d_{\text{model}}\times d_{\text{model}}}, \qquad \mathrm{rank}(W_Q W_K^\top) \le d_{\text{head}},
\label{eq:low_rank_lift}
\end{equation}
which is rank-deficient and therefore not an element of $\mathrm{GL}(d_{\text{model}})$. If $M_h$ is read gauge-theoretically, its bilinear content is that of the pair-dependent untied-route object $M_{ij} = \Omega_{ij}^{-\top}\Sigma_j^{-1}$ of Eq.~\eqref{eq:gauge_bilinear}, which reduces to $\sigma^{-2}\Omega_{ij}^{-\top}$ under isotropic covariances; within the trivial-frame route of this subsection $M_h$ is instead a separately learned head-space object (hypothesis (R3)), and its only gauge content is the block-diagonal restriction it inherits. The isometric columns of $U_Q, U_K$ select the per-head image subspaces and carry no gauge content. The gauge group of the multi-head architecture is therefore the block-diagonal subgroup $\mathrm{GL}(d_{\text{head}})^H \subset \mathrm{GL}(d_{\text{model}})$ identified in Eq.~\eqref{eq:multihead_gauge_group}.

A factorization of the form \eqref{eq:head_space_kernel_participatory} always exists for any invertible $M_h$: every element of $\mathrm{GL}(d_{\text{head}})$ admits a decomposition $A B^\top$ with $A, B \in \mathrm{GL}(d_{\text{head}})$ (e.g., via the square SVD $M_h = U\Lambda V^\top = (U\Lambda^{1/2})(V\Lambda^{1/2})^\top$, where invertibility of $M_h$ guarantees positive singular values, so $A$ and $B$ are themselves invertible). The factorization is unique only up to the right action $A_Q \to A_Q G$, $A_K \to A_K G^{-\top}$ for $G \in \mathrm{GL}(d_{\text{head}})$; imposing the balanced normalization $A_Q^\top A_Q = A_K^\top A_K$ (e.g., the SVD form $A_Q = U\Lambda^{1/2}$, $A_K = V\Lambda^{1/2}$) reduces this freedom to a common right orthogonal action $A_Q \to A_Q R$, $A_K \to A_K R$ with $R \in \mathrm{O}(d_{\text{head}})$. Directions in this stabilizer leave $M_h$, and hence the attention computation, unchanged, so training cannot select among them except through the implicit bias of gradient descent toward balanced factorizations: gradient flow on factorized objectives conserves $A_Q^\top A_Q - A_K^\top A_K$, so balanced initializations remain balanced~\cite{du2018balanced, arora2018implicit}.

Define query and key vectors as
\begin{equation}
Q_i \equiv \mu_i^\top W_Q \in \mathbb{R}^{1\times d_k}, \qquad K_j \equiv \mu_j^\top W_K \in \mathbb{R}^{1\times d_k},
\end{equation}
where $W_Q$ acts on the right of the row vector $\mu_i^\top$; the untied route of Section~\ref{sec:untied_qk} wrote the equivalent left action $Q_i = W_Q^{(i)}\mu_i$. Then
\begin{equation}
Q_i K_j^\top = \mu_i^\top W_Q W_K^\top \mu_j = \frac{1}{\sigma^2}\mu_i^\top M \mu_j.
\end{equation}
Thus the logit term reduces to
\begin{equation}
\tilde{s}_{ij} = Q_i K_j^\top - \frac{1}{2\sigma^2}\|\mu_j\|^2.
\end{equation}
Under hypothesis (N) in its exact pre-projection form, $\|\mu_j\|$ is constant in $j$, the second term is row-constant, and it cancels in the softmax (the first regime above); outside the regimes of (N) it survives as a key-dependent bias of the same order as the content logits.

\paragraph{Temperature scaling.}
Since $\sigma^{-2}$ is absorbed into the learned projections $W_Q, W_K$, the effective logit $\tilde{s}_{ij} = Q_i K_j^\top$ carries no residual dependence on $\sigma$. The softmax $\beta_{ij} = \mathrm{softmax}_j(\tilde{s}_{ij}/\tau)$ must therefore set $\tau$ to account only for the dimensional scaling of the dot product.

In high-dimensional spaces, dot products $Q_iK_j^\top$ have standard deviation $O(\sqrt{d_k})$ (for unit-variance components), while key-bias fluctuations are $O(1)$ after layer normalization. Normalizing pre-softmax logits to $O(1)$ variance therefore requires the dimensional factor $\sqrt{d_k}$. In the canonical form of the manuscript the temperature is factorized as $\tau = \kappa\sqrt{d_k}$ with $\kappa > 0$ a learnable scalar (the canonical $\tau = \kappa\sqrt{K_q}$ of Section~\ref{sec:complete_free_energy} and Eq.~\eqref{eq:pointwise_free_energy}, specialized here via the per-head identification $K_q \to d_{\mathrm{head}} = d_k$ of this transformer reduction); the $\sqrt{d_k}$ factor is the dimensional contribution recovered from the variance argument above, and $\kappa$ controls the residual sharpness that the network adapts during training. Combining all simplifications:
\begin{equation}
\boxed{\beta_{ij} = \mathrm{softmax}_j\left(\frac{Q_i K_j^\top}{\kappa\sqrt{d_k}}\right)}\quad\text{(canonical form; }\kappa=1\text{ recovers~\cite{vaswani2017attention}).}
\end{equation}
The recovery of standard scaled dot-product attention is therefore exact under (R1)--(R3) and (N) up to the learnable temperature $\kappa$, proving Theorem~\ref{thm:transformer_recovery}.

\subsection{Value Aggregation}\label{sec:value_aggregation_itfb}

The mixture-of-sources generative model (Section~\ref{sec:mixture_derivation}) specifies that the component distribution for source $j$ in agent $i$'s frame has mean $\Omega_{ij}\mu_j$. The posterior mean under mixture responsibilities $\beta_{ij}$ is therefore
\begin{equation}
\hat\mu_i = \sum_j \beta_{ij} \Omega_{ij}\mu_j,
\label{eq:glk_mixture_aggregation_itfb}
\end{equation}
which is the standard mixture-of-Gaussians expectation. Thus the general gauge-covariant value object is pair-specific:
\begin{equation}
V_{ij} := \Omega_{ij}\mu_j,
\qquad
\hat\mu_i = \sum_j \beta_{ij}V_{ij}.
\end{equation}
This is the exact mixture-of-Gaussians expectation implied by the consensus construction. It is not, in general, representable as a standard source-only value projection $V_j=W_V^\top\mu_j$, because $\Omega_{ij}=U_iU_j^{-1}$ depends on the receiving agent $i$. The standard transformer value update is recovered only under the shared-frame limit $\Omega_{ij}=I$ of hypothesis (R2), together with a separately introduced source-only value map $W_V$ (fence F1 below).

\subsection{Complete Attention Formula}

Under the assumptions above, the gauge-theoretic message-passing
\begin{equation}
\beta_{ij} \propto \exp\left(-\frac{1}{\tau}D_{\mathrm{KL}}(q_i \| \Omega_{ij} q_j)\right), \quad \hat\mu_i = \sum_j \beta_{ij} \Omega_{ij}\mu_j,
\end{equation}
reduces to
\begin{equation}
\boxed{\mathrm{Attention}(Q,K,V) = \mathrm{softmax}\left(\frac{QK^\top}{\kappa\sqrt{d_k}}\right)V},
\end{equation}
where $V$ stacks the transported means $\Omega_{ij}\mu_j = \mu_j$, the identity value map left by the shared-frame limit; composing with the separately introduced source-only value map $W_V$ of the preceding subsection, and setting $\kappa = 1$, recovers the standard scaled dot-product attention of~\cite{vaswani2017attention}. This is the canonical dot-product attention mechanism in transformers. To be precise about what is gauge-derived (this is the full statement of fence F1; Section~\ref{sec:epistemic_status}): the score $QK^\top$ and its softmax weighting follow from the KL-consensus stationary point under the hypotheses (R1)--(R3) and (N) of Theorem~\ref{thm:transformer_recovery}, whereas the value map $V$ (via $W_V$) and the post-attention output mixing $W_O$ retained in the sweep configuration are standard-transformer objects re-introduced here, not consequences of the gauge transport. The boxed identity is therefore a consistency embedding of vanilla attention into the KL-attention family at the score-and-weighting level, a recovery in a limit rather than a derivation of attention from gauge theory alone, and not a parameter-free reconstruction of the full $\mathrm{Attention}(Q,K,V)$ operator.

The limits are deliberately aggressive: (R1) collapses the statistical manifold into a basic Euclidean space, and (R3) fills the cross-coupling slot vacated by the trivial transport of (R2) with a separately learned bilinear (only the scale $\sigma^{-2}$ is absorbed into $W_Q, W_K$). Each limit discards a specific geometric degree of freedom and each can be independently relaxed to produce a generalization of standard attention. The gauge-theoretic framework provides a taxonomy of design choices, with standard attention as the three-limit simplification and richer architectures (non-isotropic beliefs, learned gauge transport, position-dependent connections) occupying intermediate points.

\paragraph{Multi-head attention as a block-diagonal gauge restriction.}
The multi-head attention of~\cite{vaswani2017attention} partitions the embedding dimension into $H$ heads of dimension $d_{\text{head}} = d_{\text{model}}/H$, each computing attention independently with its own projection matrices and concatenating the results. Each head $a \in \{1,\ldots,H\}$ carries an independent $\mathrm{GL}(d_{\text{head}})$ gauge transformation $\Omega^{(a)}$ acting on its $d_{\text{head}}$-dimensional subspace, with no cross-head mixing under transport; assembling the per-head transformations into a block-diagonal element of $\mathrm{GL}(d_{\text{model}})$ identifies the multi-head gauge group as the direct product
\begin{equation}
G_{\text{multi-head}} = \mathrm{GL}(d_{\text{head}})^H \subset \mathrm{GL}(d_{\text{model}}),
\label{eq:multihead_gauge_group}
\end{equation}
a subgroup of dimension $H \cdot d_{\text{head}}^2$ inside the $d_{\text{model}}^2$-dimensional ambient group, consistent with the low-rank lift of Eq.~\eqref{eq:low_rank_lift}. The corresponding Lie-algebra restriction is $\bigoplus_a \mathfrak{gl}(d_{\text{head}}) \subset \mathfrak{gl}(d_{\text{model}})$, and the off-diagonal complement $\mathfrak{m}$ of dimension $d_{\text{model}}^2 - H \cdot d_{\text{head}}^2$ in the decomposition $\mathfrak{gl}(d_{\text{model}}) = \bigoplus_a \mathfrak{gl}(d_{\text{head}}) \oplus \mathfrak{m}$ is the cross-head mixing sector that the multi-head architectural choice projects out of the attention computation (the standard output projection $W_O$ mixes head outputs \emph{after} attention rather than restoring this complement before it).

\paragraph{Rotary position embeddings as position-dependent gauge frames.}
Rotary position embeddings (RoPE) \cite{su2024roformer} apply position-dependent rotations to the query and key vectors,
\begin{equation}
Q_i \to R(\theta_i) Q_i, \qquad K_j \to R(\theta_j) K_j,
\label{eq:rope_apply}
\end{equation}
with $R(\theta_m) \in \mathrm{SO}(2)^{d_k/2}$ a block-diagonal rotation through angles $m\theta_n$ at frequencies $\theta_n = 10000^{-2n/d_k}$. The logit kernel becomes $Q_i^\top R(\theta_i)^\top R(\theta_j) K_j = Q_i^\top R(\theta_{j-i}) K_j$, depending only on the relative position $j - i$. Identifying $\phi_i^{(\mathrm{pos})} \in \mathfrak{so}(2)^{d_k/2} \subset \mathfrak{gl}(d_k)$ as a position-dependent gauge frame coefficient with $R(\theta_i) = \exp(\phi_i^{(\mathrm{pos})})$, the relative rotation $R(\theta_{j-i})$ has the gauge-transport form $\exp(-\phi_i^{(\mathrm{pos})})\exp(\phi_j^{(\mathrm{pos})}) = \Omega_{ij}^{-1}$ (equivalently $\Omega_{ij}^\top$ for the orthogonal blocks), so RoPE is the position-dependent gauge specialization of the framework's transport $\Omega_{ij} = U_iU_j^{-1}$ restricted to the compact subgroup $\mathrm{SO}(2)^{d_k/2}$. This is the trivial-frame-route counterpart of the untied-route identification with rotary structure in Section~\ref{sec:untied_qk}, with $U_i = R(\theta_i) = \exp(\phi_i^{(\mathrm{pos})})$; there the closure condition on the covariance field carried the correspondence, here the abelian commutativity of the $\mathfrak{so}(2)$ blocks does. The trivial-frame limit $U_i = U$ removes the position dependence and recovers the constant-gauge reduction of this subsection; the full $\mathrm{GL}(d_k)$ transport allows the gauge to be content- as well as position-dependent and learned. The collapse $\exp(-\phi_i^{(\mathrm{pos})})\exp(\phi_j^{(\mathrm{pos})}) = \exp(\phi_j^{(\mathrm{pos})} - \phi_i^{(\mathrm{pos})})$ that produces the relative-position-only kernel $R(\theta_{j-i})$ holds because the $\mathfrak{so}(2)$ blocks commute; the relative-position-only property is therefore specific to the abelian rotation sector. A non-abelian $\mathrm{GL}(d_k)$ lift retains the position- and content-dependent transport but loses the exact $f(i-j)$ dependence, since the Baker-Campbell-Hausdorff term $-\tfrac{1}{2}[\phi_i^{(\mathrm{pos})}, \phi_j^{(\mathrm{pos})}]$ no longer vanishes. RoPE is thus recovered exactly on the abelian rotation sector; the non-abelian $\mathrm{GL}(d_k)$ lift generalizes the transport but forfeits the relative-position-only form.

\paragraph{GLU-family gating in the VFE mean update.}
The per-edge contribution of neighbor $j$ to the receiver mean update under the entropy-suppressed surrogate has the form
\begin{equation}
\Delta\mu_i\big|_j \propto \underbrace{e_{ij}}_{\text{linear residual}} \cdot \underbrace{\frac{\exp(-\|e_{ij}\|^2_{\Sigma_j^{-1}} / 2\tau)}{\sum_k \exp(-\|e_{ik}\|^2_{\Sigma_k^{-1}}/2\tau)}}_{\text{Boltzmann gate over KL energies}},
\label{eq:vfe_glu}
\end{equation}
with $e_{ij} = \mu_i - \Omega_{ij}\mu_j$ the transported residual. This is a Boltzmann-gated linear unit: a linear function of the prediction error gated by a softmax of KL energies computed from the same input. In the binary-source limit ($n_s = 2$) the Boltzmann gate collapses to the logistic sigmoid and the per-edge map becomes $\Delta\mu_i\big|_j \propto e_{ij}\cdot\sigma\big((\|e_{ik}\|^2_{\Sigma_k^{-1}} - \|e_{ij}\|^2_{\Sigma_j^{-1}})/2\tau\big)$: an odd, redescending (residual-suppressing) gated map whose gate argument is a difference of quadratic energies in the residual, not the residual itself. This is GLU-style gating, in which value and gate are different functions of the same input; it is structurally analogous to, but not equal to, the SiLU / Swish form $f(x) = x\cdot\sigma(x)$ used in transformer feed-forward sublayers, whose gate is linear in its argument, and to the GELU $f(x) = x\cdot\Phi(x)$, which draws its gate from the same $x\cdot\mathrm{CDF}(\cdot)$ gating family. The framework motivates gated activations of this family (they implement information-geometric gating of prediction errors) without recovering SiLU or GELU exactly; iterated VFE updates with re-computed $\beta$ occupy the same block position as FFN depth, but we do not claim the iterated map matches the function class of an MLP sublayer, and no exact activation is recovered. Under the canonical entropy-augmented $\mathcal{F}$ at Eq.~\eqref{eq:free_energy_functional_final} the explicit softmax-gradient correction $\partial\beta_{ij}/\partial\mu_i$ cancels against the entropy term by the row-Lagrangian KKT identity at $\beta = \beta^*$; the autograd-on-surrogate route retains it as the additional centered-softmax nonlinearity of Section~\ref{sec:envelope_theorem}.

\paragraph{Generalization to R\'enyi $\alpha$-divergence.}
The construction uses the KL divergence as the alignment functional. The framework extends, on the domain where the divergence is finite, to the one-parameter R\'enyi family $D_\alpha(q \| p)$ with order $\alpha \in (0, \infty) \setminus \{1\}$, and the $\alpha \to 1$ limit recovers the KL. For independent diagonal Gaussians the R\'enyi divergence admits the closed form
\begin{equation}
D_\alpha(q \| p) = \frac{1}{2}\left[ \alpha \sum_k \frac{(\mu_p^k - \mu_q^k)^2}{\tilde\sigma_k} + \frac{1}{\alpha - 1} \sum_k \left((1-\alpha)\log \sigma_q^k + \alpha \log \sigma_p^k - \log \tilde\sigma_k\right) \right],
\label{eq:renyi_diagonal_itfb}
\end{equation}
with the blended variance $\tilde\sigma_k = (1-\alpha)\sigma_q^k + \alpha\sigma_p^k$. The closed form is valid only where the blend is positive, $\tilde\sigma_k > 0$ for every $k$: this holds automatically for $\alpha \in (0,1)$, where both coefficients are positive, but for $\alpha > 1$ the coefficient $(1-\alpha)$ is negative and positivity becomes a genuine condition on the variance pair; where it fails, $D_\alpha = +\infty$ and the formula does not apply. The full-covariance generalization follows by replacing the per-coordinate blend with the matrix blend $\tilde\Sigma = (1-\alpha)\Sigma_q + \alpha \Sigma_p$, subject to the analogous requirement $\tilde\Sigma \succ 0$, and the transported-key version is obtained by substituting $(\mu_p, \Sigma_p) \to (\Omega_{ij}\mu_j, \Omega_{ij}\Sigma_j\Omega_{ij}^\top)$. The attention construction
\begin{equation}
\beta_{ij}^{(\alpha)} = \mathrm{softmax}_j\left(-\frac{D_\alpha(q_i \| \Omega_{ij} q_j)}{\tau}\right)
\label{eq:renyi_attention_itfb}
\end{equation}
is unchanged in form; the KL row $\alpha = 1$ is the standard limit and other values select different mass/coverage trade-offs (the choice $\alpha = 1/2$ gives the Bhattacharyya distance $D_{1/2} = -2\log\int\sqrt{q_i \Omega_{ij}q_j}$, i.e.\ minus twice the log of the Bhattacharyya coefficient). On the finiteness domain just described, the gradient derivations in this section and in the appendix carry over by replacing $\partial D_{\mathrm{KL}}/\partial \theta$ with $\partial D_\alpha/\partial \theta$; specifically the product-rule chain Eq.~\eqref{eq:alpha_product_rule_itfb} for the state-dependent prior coupling retains its form with $D_{\mathrm{KL}} \to D_\alpha$. One element does not transfer automatically: the mixture-model responsibility derivation of value aggregation (Section~\ref{sec:value_aggregation_itfb}) is specific to the KL divergence, so under $D_\alpha$ the aggregation rule is retained as a definition rather than rederived from a variational mixture argument. The order parameter $\alpha$ here is distinct from the per-agent prior precision $\alpha_i$ of Section~\ref{sec:state_dependent_precision}; the subscript disambiguates.

This completes the descent to the bottom rung: the general gauge architecture on a point, trained as a working language model. The next section reports the empirical scaling behavior of that architecture; the sections after it climb back up, first characterizing the second-variation stiffness of the single-scale functional and then constructing the coarse-graining that stacks such rungs into a tower.
\section{Results: Aggregate Dimension Sweep on WikiText-103}\label{sec:scaling_validation}

The bundled aggregate table characterizes the gauge architecture across eleven embedding dimensions $K_q\in\{10,20,30,40,50,60,70,80,90,100,120\}$. It contains 32 runs in total: two at $K_q=90$ and three at every other dimension. Each marker in Fig.~\ref{fig:scaling_main} is an aggregate test-perplexity mean and each bar is one standard deviation across the available runs at that dimension. The table reports a constant mean token count of $122{,}880{,}000$ across dimensions, so the sweep is descriptive and iso-token at the aggregate level. Raw per-seed gauge outputs and run configurations are unavailable; the aggregate table cannot support resampling inference, seed-identity claims, or reconstruction of the training procedure.

Nonlinear least squares on the eleven aggregate means, using the bounded point-fit model $\mathrm{PPL}=aK_q^b+c$, gives
\begin{align*}
a &= 1805.5529, \\
b &= -1.048919, \\
c &= 61.1699,
\end{align*}
with $R^2=0.99981995$. This coefficient of determination describes the eleven aggregate means under the chosen three-parameter family; it is not an uncertainty measure or evidence that the family is universal. The point estimate is close to an inverse-$K_q$ dependence over this grid, but the fitted floor and exponent are coupled and the raw runs needed for a sampling statement are absent. The result is therefore a within-architecture descriptive trend, not a scaling-law universality claim.

\begin{figure}[t]
\centering
\includegraphics[width=\linewidth]{figures/pifb2_scaling.pdf}
\caption{\textbf{Aggregate WikiText-103 dimension sweep.} Markers are aggregate test-perplexity means and bars are one standard deviation across the available runs. The $K_q=90$ point contains two runs; every other point contains three, for 32 runs total. The curve is the aggregate-level point fit $\mathrm{PPL}=1805.5529K_q^{-1.048919}+61.1699$. It is descriptive and carries no bootstrap band or sampling interval.}
\label{fig:scaling_main}
\end{figure}

At $K_q=120$ the aggregate mean test perplexity is approximately $72.71$. No comparison with published word-level WikiText-103 perplexities is made because this sweep uses a different prediction unit; tokenizer-compatible comparison requires a new controlled experiment with identical preprocessing, split, optimization budget, and evaluation code. The fitted floor is likewise a parameter of this finite, aggregate iso-token sweep rather than an architectural ceiling.

